\documentclass[12pt,a4paper]{ltjarticle}


% ------------------------------
% 数学パッケージ
% ------------------------------
\usepackage{amsmath,amssymb,amsthm,amsfonts}
\usepackage{bm} % 太字ベクトル
\usepackage{mathrsfs} % 花文字
\usepackage{mathtools}  % amsmath の拡張
\usepackage{physics}  % 量子力学向け記法（bra-ket など）

% ------------------------------
% レイアウト
% ------------------------------
\usepackage{geometry}
\geometry{margin=25mm}

\usepackage{setspace}
\setstretch{1.3}  % 行間を広げて読みやすく

% ------------------------------
% 見出しデザイン
% ------------------------------
\usepackage{titlesec}
\titleformat{\section}{\Large\bfseries}{\thesection}{1em}{}
\titleformat{\subsection}{\large\bfseries}{\thesubsection}{1em}{}
\titleformat{\subsubsection}{\normalsize\bfseries}{\thesubsubsection}{1em}{}

% ------------------------------
% 定理環境
% ------------------------------
\newtheorem*{theorem}{定理}
\newtheorem*{lemma}{補題}
\newtheorem*{proposition}{命題}
\newtheorem*{corollary}{系}

\theoremstyle{definition}
\newtheorem*{definition}{定義}
\newtheorem*{remark}{注意}
\newtheorem*{example}{例}

% ------------------------------
% ハイパーリンク
% ------------------------------
\usepackage[hidelinks]{hyperref}

% ------------------------------
% 追加の便利パッケージ
% ------------------------------
\usepackage{enumitem}% 箇条書きの制御
\usepackage{cite} % 文献引用の整形
\usepackage{url}  % URL の整形
\begin{document}

% ------------------------------
% タイトル
% ------------------------------
\title{群環・複素数体上の線形結合の許容と作用素的構造の拡張}
\author{御館 凛太朗\\\scriptsize{ontachi@outlook.jp}}
\date{2026年3月1日\\\scriptsize{DOI:10.5281/zenodo.19186600}}
\maketitle
\begin{abstract}
  本稿では、群環 $\mathbb{C}G$ における線形結合の構造を、代数学、解析学、および量子力学の観点から検討する。群環の形式和が、左正則表現を通じてヒルベルト空間 $\ell^2(G)$ 上の有界作用素として機能することを示し、その構造を既約群 $C^*$-環および von Neumann 環へと拡張することで、スペクトル性質やトレース構造を論じる。さらに、群環の畳み込み積を量子遷移振幅の離散的な経路積分として解釈する視点を提示する。本研究は、非可換調和解析と量子作用素論の境界を埋めるものであり、元のスペクトル半径が群の幾何学的・確率論的特性を反映していることを明らかにする。\\\scriptsize{© 2026 Rintaro Ontachi. This work is licensed under a Creative Commons Attribution-NonCommercial-NoDerivs 4.0 International License.}
\end{abstract}
\pagebreak


\section{序論：線形結合の現代的意義}

群環 $\mathbb{C}G$ の元は形式的な有限和
\[
x = \sum_{g \in G} a_g g
\]
として表されるが、この単純な形の背後には、非可換性・作用素性・量子性が表れている。特に、線形結合という操作は、単なる代数的操作ではなく、量子力学における「状態の重ね合わせ」や「遷移振幅の干渉」と同型の構造を持つ。

この形式和は一見すると単なる記号操作に見えるが、実際には群の代数構造と複素線形空間の構造の表現である。群の元 $g$ は基底ベクトルとなり、複素係数 $a_g$ はその基底方向の“重み”を表す。したがって、群環の元は群の元を基底とする複素ベクトル空間であり、同時に群の積に基づく非可換代数でもある。

\begin{definition}[群環の形式和]
群 $G$ に対し、有限個の $a_g\in\mathbb{C}$ が非零となる形式和
\[
x=\sum_{g\in G}a_g g
\]
を群環 $\mathbb{C}G$ の元と呼ぶ。
\end{definition}

この定義は単純であるが、以下の三つの重要な意味を同時に含んでいる。

\begin{itemize}
\item 群の元が線形独立な基底として振る舞うため、$\mathbb{C}G$ は巨大な複素ベクトル空間である。
\item 群の積が畳み込みとして形式和に作用し、非可換代数構造を与える。
\item 複素係数が量子力学的干渉（位相の重ね合わせ）を可能にする。
\end{itemize}

線形結合の係数が複素数であることは、位相の干渉を許容に対して良い性質を示す。実数係数では表現できない複素位相の回転が、群環の元の積において有用である。例えば、二つの元
\[
x = \sum_{g} a_g g, \qquad y = \sum_{h} b_h h
\]
の積は
\[
xy = \sum_{k \in G} \left( \sum_{gh = k} a_g b_h \right) k
\]
で与えられるが、この内側の和は量子力学における経路和に極めて近い構造を持つ。すなわち、$k$ に到達するすべての経路 $g \to h$ の振幅が複素数として干渉し合う。

\begin{remark}[経路和との同型性]
量子力学における遷移振幅
\[
A(g\to k)=\sum_{h}A(g\to h)A(h\to k)
\]
と群環の積における係数
\[
c_k=\sum_{gh=k}a_g b_h
\]
は完全に同型である。したがって群環の積は、量子遷移の離散モデルとみなすことができる。
\end{remark}

このように、群環の線形結合は、単なる代数的操作ではなく、量子系の遷移振幅の重ね合わせと同型の構造を持つ。したがって、群環を作用素として扱うことは、量子力学的な視点からも自然であり、現代数学の多くの分野と深く結びついている。

さらに、群環は確率論的にも自然な意味を持つ。係数 $a_g$ を確率振幅とみなすと、群環の積はランダムウォークの遷移確率の合成と一致する。この観点は、後に非可換確率論や自由確率論との接続を与える。

\begin{example}[ランダムウォークとの関係]
確率測度 $\mu(g)$ を係数とする元
\[
x=\sum_{g}\mu(g)g
\]
の $n$ 乗は $n$ ステップ後の遷移確率を与える：
\[
x^n=\sum_{g}\mu^{*n}(g)g.
\]
\end{example}

本稿では、群環を単なる代数的対象としてではなく、ヒルベルト空間上の作用素として再構成し、その解析的性質を徹底的に明らかにする。これにより、群の対称性が作用素のスペクトル構造にどのように反映されるかを精密に記述する。

\begin{quote}
群環は「代数」「解析」「量子」の三つの視点が交差する場であり、本稿はその交点を体系的に記述することを目的とする。
\end{quote}

%%%%%%%%%%%%%%%%%%%%%%%%%%%%%%%%%%%%%%%%%%%%%%%%%%%%%%%%%%%%%%
\section{群環の代数的・ベクトル空間的構造}

\subsection{群環の元の構造と畳み込み積}

群環 $\mathbb{C}G$ の元 $x$ は、係数関数 $a: G \to \mathbb{C}$ によって特徴付けられる。すなわち、
\[
x = \sum_{g \in G} a_g g
\]
であり、$a_g$ が 0 でない $g$ は有限個である。

この形式和は、単に「有限個の項を持つ複素線形結合」というだけではなく、群の代数構造と線形空間構造が緊密に融合した高度な数学的対象である。特に、群の元が線形独立な基底として振る舞うため、$\mathbb{C}G$ は巨大な無限次元ベクトル空間となる。

\begin{definition}[有限支援関数としての群環]
群環 $\mathbb{C}G$ は、有限支援関数
\[
a:G\to\mathbb{C},\qquad \mathrm{supp}(a)=\{g\in G\mid a(g)\neq 0\}\ \text{有限}
\]
全体と同一視できる。
\end{definition}

この視点は、後に $\ell^2(G)$ 上の畳み込み作用として群環を実現する際に本質的な役割を果たす。

---

積構造は群の積に基づく畳み込みであり、
\[
(xy)(k) = \sum_{gh=k} a_g b_h
\]
と書ける。この畳み込みは、関数空間 $L^1(G)$ における畳み込み
\[
(f * g)(x) = \int_G f(y) g(y^{-1}x) \, dy
\]
の離散版であり、非可換性がそのまま代数構造に反映される。

\begin{proposition}[畳み込み積の結合性]
任意の $x,y,z\in\mathbb{C}G$ に対して
\[
(xy)z = x(yz)
\]
が成立する。
\end{proposition}

\begin{proof}
両辺の $k$ 成分を比較すると
\[
((xy)z)(k)=\sum_{uv=k}(xy)(u)z(v)=\sum_{uv=k}\left(\sum_{gh=u}a_g b_h\right)c_v,
\]
\[
(x(yz))(k)=\sum_{gh=k}a_g(yz)(h)=\sum_{gh=k}a_g\left(\sum_{uv=h}b_u c_v\right).
\]
群の積の結合性 $g(uv)=(gh)v$ により一致する。
\end{proof}

特に、$G$ が非可換群である場合、$xy \neq yx$ が一般に成立する。これは、群環が非可換環の典型例であることを示している。

\begin{example}[非可換性の具体例]
$G=S_3$ とすると、$(12)(23)\neq(23)(12)$ であるため、
\[
(12)+(23) \quad\text{の積は順序に依存する。}
\]
\end{example}

---

\subsection{線形結合の意味の深化}

線形結合
\[
x = \sum_{g \in G} a_g g
\]
は、単に「形式的な和」ではなく、群の元を基底とする複素ベクトル空間としての構造を明確に示している。

ここで重要なのは、群の元 $g$ が「互いに独立な基底ベクトル」として振る舞う点である。すなわち、
\[
\sum_{g \in G} a_g g = 0
\]
が成立するならば、すべての $a_g = 0$ が従う。これは、群の元が線形独立であることを意味する。

\begin{proposition}[群の元の線形独立性]
$\{g\mid g\in G\}$ は $\mathbb{C}G$ の基底である。
\end{proposition}

\begin{proof}
$\sum_g a_g g=0$ と仮定し、$\ell^2(G)$ の基底 $e_h$ に作用させると
\[
0=\left(\sum_g a_g g\right)e_e=\sum_g a_g e_g.
\]
$\{e_g\}$ の線形独立性より $a_g=0$ が従う。
\end{proof}

この線形独立性は、後に導入する正則表現や Fourier 解析において極めて重要な役割を果たす。特に、群環の Fourier 変換は、基底 $\{g\}$ の線形独立性を前提として行列ブロックへの分解を可能にする。

---

\subsection{対合の導入と代数構造}

群環の対合は
\[
x^* = \sum_{g \in G} \overline{a_g} g^{-1}
\]
で定義される。この対合は、以下の性質を満たす。
\[
(x^*)^* = x, \qquad (xy)^* = y^* x^*
\]
これにより $\mathbb{C}G$ は $^*$-代数となり、作用素環への自然な拡張が可能となる。

\begin{lemma}[対合の反転性]
任意の $g\in G$ に対して $(g^{-1})^{-1}=g$ であるため、
\[
(x^*)^*=\sum_g a_g g=x.
\]
\end{lemma}

\begin{proposition}[対合と畳み込みの関係]
\[
(xy)^*=y^*x^*
\]
が成立する。
\end{proposition}

\begin{proof}
\[
(xy)^*=\sum_k \overline{(xy)(k)}k^{-1}
=\sum_{gh=k}\overline{a_g b_h}k^{-1}
=\sum_{g,h}\overline{b_h}h^{-1}\,\overline{a_g}g^{-1}
=y^*x^*.
\]
\end{proof}

特に、$x = x^*$ が成立する場合、$x$ は自己共役であり、後に述べるスペクトル理論において「実スペクトル」を持つことが保証される。

\begin{remark}[自己共役元の物理的解釈]
自己共役元は量子力学における観測可能量（observable）に対応し、  
固有値は観測値、固有射影は測定結果の確率を与える。
\end{remark}

%%%%%%%%%%%%%%%%%%%%%%%%%%%%%%%%%%%%%%%%%%%%%%%%%%%%%%%%%%%%%%
\section{正則表現とヒルベルト空間への埋め込み}

群環 $\mathbb{C}G$ を作用素として理解するための最も基本的な方法は、
ヒルベルト空間 $\ell^2(G)$ 上の左正則表現を用いることである。
この構成は単に「行列表示を与える」というだけではなく、
群環の解析的性質を引き出すための本質的な道具となる。

左正則表現は、群の抽象的な構造を具体的な作用素として実現するための最も自然な方法であり、
群環の作用素論的研究の基盤を形成する。
特に、群 von Neumann 環 $L(G)$ や reduced group $C^*$-algebra $C_r^*(G)$ の構成は、
この正則表現を通じて初めて可能となる。

%%%%%%%%%%%%%%%%%%%%%%%%%%%%%%%%%%%%%%%%%%%%%%%%%%%%%%%%%%%%%%
\subsection{ヒルベルト空間の構造}

まず、ヒルベルト空間 $\ell^2(G)$ を
\[
\ell^2(G) = \left\{ f : G \to \mathbb{C} \,\middle|\,
\sum_{g \in G} |f(g)|^2 < \infty \right\}
\]
と定義する。
内積は
\[
\langle f, h \rangle = \sum_{g \in G} f(g)\overline{h(g)}
\]
で与えられる。

この空間は、群の元を添字とする無限次元ヒルベルト空間であり、
群の構造を解析的に扱うための自然な舞台となる。
特に、群の元に対応する標準基底 $\{e_g\}_{g\in G}$ は、
群の構造をそのままヒルベルト空間に持ち込む役割を果たす。

\begin{definition}[標準基底]
各 $g\in G$ に対し
\[
e_g(h)=
\begin{cases}
1 & (h=g),\\
0 & (h\neq g)
\end{cases}
\]
で定義される関数 $e_g$ を $\ell^2(G)$ の標準基底と呼ぶ。
\end{definition}

\begin{lemma}[正規直交性]
標準基底は
\[
\langle e_g, e_h\rangle=\delta_{g,h}
\]
を満たす正規直交系である。
\end{lemma}

\begin{proof}
定義より明らか。
\end{proof}

さらに、任意の $f\in\ell^2(G)$ は
\[
f=\sum_{g\in G}f(g)e_g
\]
と展開できるため、標準基底は完全系である。

\begin{remark}[群の構造と $\ell^2(G)$]
群の元が添字として現れるため、$\ell^2(G)$ は群の作用を自然に受け入れる。
この性質が、後に述べる正則表現の構成を可能にする。
\end{remark}

%%%%%%%%%%%%%%%%%%%%%%%%%%%%%%%%%%%%%%%%%%%%%%%%%%%%%%%%%%%%%%
\subsection{左正則表現の定義と基本性質}

左正則表現 $\lambda : G \to \mathcal{B}(\ell^2(G))$ は
\[
(\lambda(g)f)(h) = f(g^{-1}h)
\]
で定義される。
これは「左からの平行移動」を表す作用であり、
群の構造をそのままヒルベルト空間上の作用として実現している。

特に、
\[
\lambda(g)e_h = e_{gh}
\]
が成立するため、基底の添字が群の積に従って移動する。

\begin{proposition}[群準同型性]
任意の $g,h\in G$ に対して
\[
\lambda(gh)=\lambda(g)\lambda(h)
\]
が成立する。
\end{proposition}

\begin{proof}
$f\in\ell^2(G)$ に対し、
\[
(\lambda(gh)f)(x)=f((gh)^{-1}x)=f(h^{-1}g^{-1}x)
\]
一方、
\[
(\lambda(g)\lambda(h)f)(x)=\lambda(h)f(g^{-1}x)=f(h^{-1}(g^{-1}x))
\]
であるため一致する。
\end{proof}

この作用はユニタリ作用素であり、
\[
\lambda(g)^* = \lambda(g^{-1})
\]
が成立する。

\begin{lemma}[ユニタリ性]
\[
\lambda(g)^*\lambda(g)=\lambda(g)\lambda(g)^*=I
\]
が成立する。
\end{lemma}

\begin{proof}
内積の保存性
\[
\langle \lambda(g)f,\lambda(g)h\rangle=\langle f,h\rangle
\]
より従う。
\end{proof}

\begin{remark}[置換行列としての理解]
$\lambda(g)$ は無限次元の置換行列であり、行列の各行・列にちょうど 1 個の $1$ が現れる。
したがってユニタリである。
\end{remark}

%%%%%%%%%%%%%%%%%%%%%%%%%%%%%%%%%%%%%%%%%%%%%%%%%%%%%%%%%%%%%%
\subsection{群環の元の作用素としての実現}

一般の群環の元
\[
x = \sum_{g \in G} a_g g
\]
に対して、
\[
\lambda(x) = \sum_{g \in G} a_g \lambda(g)
\]
と定義することで、$x$ は $\ell^2(G)$ 上の有界作用素として実現される。

この作用は具体的には
\[
(\lambda(x)f)(h) = \sum_{g \in G} a_g f(g^{-1}h)
\]
で与えられる。
これは畳み込み作用
\[
(\lambda(x)f)(h) = (a * f)(h)
\]
と一致する。

\begin{definition}[畳み込み作用素]
有限支援関数 $a:G\to\mathbb{C}$ に対し、
\[
T_a f(h)=\sum_{g}a(g)f(g^{-1}h)
\]
で定義される作用素 $T_a$ を畳み込み作用素と呼ぶ。
\end{definition}

\begin{proposition}[有界性]
\[
\|\lambda(x)\|_{\mathrm{op}}\le \sum_{g}|a_g|
\]
が成立する。
\end{proposition}

\begin{proof}
Cauchy–Schwarz の不等式より
\[
|(\lambda(x)f)(h)|\le \sum_{g}|a_g||f(g^{-1}h)|
\le \left(\sum_{g}|a_g|\right)\|f\|_\infty.
\]
$\ell^2$ ノルムに対しても同様。
\end{proof}

%%%%%%%%%%%%%%%%%%%%%%%%%%%%%%%%%%%%%%%%%%%%%%%%%%%%%%%%%%%%%%
\subsection{忠実性の詳細な証明}

左正則表現が忠実であること、すなわち
\[
\lambda(x) = 0 \quad \Rightarrow \quad x = 0
\]
を示す。

実際、$\lambda(x) = 0$ と仮定すると、
特に $e_e$（単位元に対応する基底ベクトル）に作用させると
\[
\lambda(x)e_e = \sum_{g \in G} a_g e_g = 0
\]
となる。

基底 $\{e_g\}$ は線形独立であるため、
\[
a_g = 0 \quad (\forall g \in G)
\]
が従う。

したがって $x = 0$ であり、$\lambda$ は単射である。

\begin{theorem}[忠実性の強化]
$\lambda$ は $\mathbb{C}G$ のみならず、$C_r^*(G)$ にも忠実に拡張される。
\end{theorem}

\begin{proof}
$\mathbb{C}G$ 上で忠実であることは上で示した。
$C_r^*(G)$ は作用素ノルム完備化であるため、連続延長により忠実性が保存される。
\end{proof}

\begin{remark}[忠実性の意味]
忠実性は、群環の解析的性質（ノルム・スペクトル・トレース）が
正則表現を通じて正確に反映されることを保証する。
\end{remark}

%%%%%%%%%%%%%%%%%%%%%%%%%%%%%%%%%%%%%%%%%%%%%%%%%%%%%%%%%%%%%%
\section{作用素としての群環の構造}

群環 $\mathbb{C}G$ を $\ell^2(G)$ 上の作用素として実現することで、
群環は単なる代数的対象ではなく、
解析的対象としての側面を持つようになる。
この視点は、群環を $C^*$-環や von Neumann 環へと拡張する際の基盤となり、
群の幾何的性質・表現論的性質・確率論的性質が
作用素論的枠組みの中で統一的に理解される。

特に、左正則表現 $\lambda$ を通じて得られる作用素 $\lambda(x)$ は、
群環の元 $x$ を「畳み込み作用素」として実現するものであり、
その作用素ノルム・スペクトル・自己共役性などの解析的性質は
群の構造と深く結びついている。

%%%%%%%%%%%%%%%%%%%%%%%%%%%%%%%%%%%%%%%%%%%%%%%%%%%%%%%%%%%%%%
\subsection{作用素ノルムの導入}

作用素 $\lambda(x)$ のノルムは
\[
\|x\| = \|\lambda(x)\|_{\mathrm{op}}
= \sup_{\|f\|=1} \|\lambda(x)f\|
\]
で定義される。

このノルムは群環の代数的構造と深く結びついており、
特にスペクトル半径との関係
\[
r(x) = \lim_{n \to \infty} \|x^n\|^{1/n}
\]
を通じて、群の幾何的性質を反映する。

\begin{definition}[作用素ノルム]
$T\in\mathcal{B}(\ell^2(G))$ に対し
\[
\|T\|_{\mathrm{op}}=\sup_{\|f\|_2=1}\|Tf\|_2
\]
を作用素ノルムと呼ぶ。
\end{definition}

\begin{lemma}[ノルムの劣乗法性]
任意の $x,y\in\mathbb{C}G$ に対して
\[
\|xy\|\le \|x\|\|y\|
\]
が成立する。
\end{lemma}

\begin{proof}
$\lambda$ が $^*$-準同型であることと作用素ノルムの一般論より従う。
\end{proof}

\begin{remark}[群の幾何との関係]
$G$ が非可換であるほど、畳み込み作用素のノルムは複雑になり、
群の成長度・アーメン性・Kazhdan 性などの性質が
作用素ノルムに反映される。
\end{remark}

---

\subsection{自己共役元と実スペクトル}

$x = x^*$ が成立する場合、$\lambda(x)$ は自己共役作用素となり、
スペクトルは実数全体の部分集合となる。
すなわち、
\[
x = x^* \quad \Rightarrow \quad \sigma(x) \subset \mathbb{R}.
\]

これは、群環の元を「観測可能量」として解釈する際に重要である。
自己共役作用素はスペクトル分解を持ち、
その固有値は物理的には測定可能な量に対応する。

\begin{theorem}[自己共役作用素のスペクトルの実性]
自己共役作用素 $A=A^*$ のスペクトルは実数に含まれる。
\end{theorem}

\begin{proof}
$A$ が自己共役であるとき、$\langle Af,f\rangle$ は常に実数である。
一般の作用素論の結果より、自己共役作用素のスペクトルは実数に含まれる。
\end{proof}

\begin{proposition}[自己共役元のノルムとスペクトル半径]
$x=x^*$ のとき、
\[
\|x\|=\max_{\lambda\in\sigma(x)}|\lambda|.
\]
\end{proposition}

\begin{proof}
自己共役作用素に対するスペクトル定理より従う。
\end{proof}

\begin{remark}[物理的解釈]
自己共役元は量子力学における観測可能量（observable）に対応し、
固有値は観測値、固有射影は測定結果の確率を与える。
群環の自己共役元は、離散対称性を持つ量子系の Hamiltonian のモデルとして現れる。
\end{remark}

---

\subsection{スペクトル構造と群の性質}

作用素 $\lambda(x)$ のスペクトルは、群の構造と密接に関係している。
特に、ランダムウォーク作用素
\[
M=\frac{1}{|S|}\sum_{s\in S}\lambda(s)
\]
のスペクトルは、群のアーメン性・拡張性・Kazhdan 性などを反映する。

\begin{example}[ランダムウォーク作用素のスペクトル]
有限生成群 $G$ の対称生成系 $S$ に対し、
\[
M=\frac{1}{|S|}\sum_{s\in S}\lambda(s)
\]
のスペクトル半径が $1$ であるか $<1$ であるかは、
$G$ のアーメン性と一致する（Kesten の定理）。
\end{example}

\begin{theorem}[Kesten の定理（簡易版）]
$G$ がアーメンであることと、
ランダムウォーク作用素 $M$ のスペクトル半径が $1$ であることは同値である。
\end{theorem}

\begin{remark}
この結果は、群の幾何的性質が作用素のスペクトルに直接反映される典型例である。
\end{remark}

---

\subsection{群 von Neumann 環への接続}

左正則表現の像
\[
L(G)=\{\lambda(g)\mid g\in G\}''
\]
は群 von Neumann 環と呼ばれ、群の解析的研究の中心的対象となる。

\begin{definition}[群 von Neumann 環]
$L(G)$ を
\[
L(G)=\lambda(G)''
\]
と定義する。ここで $''$ は弱閉包を表す。
\end{definition}

\begin{remark}
$L(G)$ は一般には II$_1$ 型因子となり、
トレース
\[
\tau(x)=\langle x e_e, e_e\rangle
\]
を持つ。
\end{remark}

このように、群環の作用素としての構造は、
$C^*$-環・von Neumann 環・非可換確率論・調和解析など、
多くの分野への入口となる。

%%%%%%%%%%%%%%%%%%%%%%%%%%%%%%%%%%%%%%%%%%%%%%%%%%%%%%%%%%%%%%
\section{非可換 Fourier 解析への導入}

群環 $\mathbb{C}G$ を作用素として扱うためには、群の表現論を通じて
その構造を「分解」する必要がある。有限群に対しては、
既約表現の完全性により、群環の元を行列の族として理解することができる。
これが非可換 Fourier 解析の基本的アイデアである。

非可換 Fourier 解析は、可換群における古典的 Fourier 解析を一般化したものであり、
群の非可換性がそのまま「行列値 Fourier 変換」として現れる。
特に、群環の元を既約表現に沿って分解することで、
群環の積・対合・ノルム・スペクトルといった解析的性質が
行列解析の問題へと還元される。

\begin{remark}[可換群との比較]
可換群ではすべての既約表現が一次元であるため、
Fourier 変換は複素数値関数の変換となる。
非可換群では既約表現が高次元となり、
Fourier 変換は行列値関数となる。
\end{remark}

%%%%%%%%%%%%%%%%%%%%%%%%%%%%%%%%%%%%%%%%%%%%%%%%%%%%%%%%%%%%%%
\subsection{既約表現と Fourier 変換の構造}

有限群 $G$ の既約表現の集合を $\widehat{G}$ とする。
各 $\rho \in \widehat{G}$ は
\[
\rho : G \to \mathrm{GL}(d_\rho, \mathbb{C})
\]
という群準同型であり、$d_\rho$ は表現の次元である。

群環の元
\[
x = \sum_{g \in G} a_g g
\]
に対して、Fourier 変換は
\[
\widehat{x}(\rho) = \sum_{g \in G} a_g \rho(g)
\]
で定義される。

これは、群環の元を「行列のブロック対角成分」に分解する操作であり、
非可換 Fourier 解析の中心的役割を果たす。

\begin{definition}[非可換 Fourier 変換]
$x=\sum_g a_g g\in\mathbb{C}G$ に対し、
\[
\widehat{x}(\rho)=\sum_{g\in G}a_g\rho(g)
\]
を $x$ の Fourier 変換と呼ぶ。
\end{definition}

\begin{lemma}[Fourier 変換の有限性]
$a_g$ が有限支援であるため、$\widehat{x}(\rho)$ は有限和で定義され、
$d_\rho\times d_\rho$ 行列となる。
\end{lemma}

\begin{remark}[行列ブロック分解]
群環 $\mathbb{C}G$ は
\[
\mathbb{C}G \cong \bigoplus_{\rho\in\widehat{G}} M_{d_\rho}(\mathbb{C})
\]
という分解を持つ（Wedderburn の定理）。
Fourier 変換はこの同型を具体的に与える。
\end{remark}

%%%%%%%%%%%%%%%%%%%%%%%%%%%%%%%%%%%%%%%%%%%%%%%%%%%%%%%%%%%%%%
\subsection{Fourier 変換の基本性質}

Fourier 変換は以下の性質を持つ。

\paragraph{線形性}
\[
\widehat{x+y}(\rho) = \widehat{x}(\rho) + \widehat{y}(\rho)
\]

\paragraph{畳み込み積との対応}
\[
\widehat{xy}(\rho) = \widehat{x}(\rho)\,\widehat{y}(\rho)
\]
これは、群環の積が表現の行列積に対応することを意味する。

\paragraph{対合との対応}
\[
\widehat{x^*}(\rho) = \widehat{x}(\rho)^{\*}
\]

\begin{proposition}[Fourier 変換は $^*$-準同型]
\[
\widehat{xy}(\rho)=\widehat{x}(\rho)\widehat{y}(\rho),\qquad
\widehat{x^*}(\rho)=\widehat{x}(\rho)^*
\]
が成立する。
\end{proposition}

\begin{proof}
表現 $\rho$ が群準同型であることから
\[
\rho(gh)=\rho(g)\rho(h)
\]
が成立し、計算により従う。
\end{proof}

%%%%%%%%%%%%%%%%%%%%%%%%%%%%%%%%%%%%%%%%%%%%%%%%%%%%%%%%%%%%%%
\subsection{Plancherel 型等式}

Fourier 変換はノルム構造とも深く関係している。
特に、群環の $\ell^2$ ノルム
\[
\|x\|_2^2 = \sum_{g \in G} |a_g|^2
\]
は、Fourier 変換を通じて
\[
\|x\|_2^2 = \frac{1}{|G|} \sum_{\rho \in \widehat{G}} d_\rho \|\widehat{x}(\rho)\|_{\mathrm{HS}}^2
\]
と表される（HS は Hilbert–Schmidt ノルム）。

\begin{definition}[Hilbert–Schmidt ノルム]
行列 $A\in M_d(\mathbb{C})$ に対し
\[
\|A\|_{\mathrm{HS}}^2=\mathrm{Tr}(A^*A)
\]
を Hilbert–Schmidt ノルムと呼ぶ。
\end{definition}

\begin{theorem}[Plancherel 等式]
\[
\sum_{g\in G}|a_g|^2
=\frac{1}{|G|}\sum_{\rho\in\widehat{G}}d_\rho\|\widehat{x}(\rho)\|_{\mathrm{HS}}^2.
\]
\end{theorem}

\begin{proof}
有限群の正規直交性関係
\[
\sum_{g\in G}\rho(g)_{ij}\overline{\sigma(g)_{kl}}
=\frac{|G|}{d_\rho}\delta_{\rho,\sigma}\delta_{i,k}\delta_{j,l}
\]
を用いて計算する。
\end{proof}

\begin{remark}[エネルギーの分散]
Plancherel 等式は、群環の元の「エネルギー」が
既約表現の行列に分散していることを意味する。
\end{remark}

%%%%%%%%%%%%%%%%%%%%%%%%%%%%%%%%%%%%%%%%%%%%%%%%%%%%%%%%%%%%%%
%%%%%%%%%%%%%%%%%%%%%%%%%%%%%%%%%%%%%%%%%%%%%%%%%%%%%%%%%%%%%%
\section{スペクトル理論の詳細}

群環の元を作用素として扱うとき、最も重要な概念のひとつがスペクトルである。
スペクトルは作用素の「固有値の一般化」であり、作用素の解析的性質を決定する。
特に、群環の元を $\ell^2(G)$ 上の作用素として実現したとき、
そのスペクトルは群の対称性・成長度・表現論的構造と密接に関係する。

スペクトル理論は、作用素の「振る舞い」を理解するための中心的道具であり、
自己共役性・正規性・ユニタリ性といった性質がスペクトルに直接反映される。
群環の元は一般に非可換であるため、スペクトルは複素平面の複雑な領域となりうる。

%%%%%%%%%%%%%%%%%%%%%%%%%%%%%%%%%%%%%%%%%%%%%%%%%%%%%%%%%%%%%%
\subsection{スペクトルの定義と基本性質}

作用素 $\lambda(x)$ のスペクトルは
\[
\sigma(x) = \{\lambda \in \mathbb{C} \mid \lambda I - \lambda(x) \text{ が可逆でない}\}
\]
で定義される。

\begin{definition}[スペクトルと解像度集合]
作用素 $T$ に対し、
\[
\rho(T)=\{\lambda\in\mathbb{C}\mid \lambda I-T\ \text{が可逆}\}
\]
を解像度集合、$\sigma(T)=\mathbb{C}\setminus\rho(T)$ をスペクトルと呼ぶ。
\end{definition}

スペクトルは常に非空であり、コンパクト集合である。
これは一般の有界作用素に対する基本定理である。

\begin{theorem}[スペクトルの基本性質]
$T\in\mathcal{B}(H)$ に対し、
\begin{itemize}
\item $\sigma(T)$ は非空
\item $\sigma(T)$ はコンパクト
\item $\sigma(T)\subseteq \{z\in\mathbb{C}\mid |z|\le \|T\|\}$
\end{itemize}
が成立する。
\end{theorem}

特に、自己共役元 $x = x^*$ の場合、
\[
\sigma(x) \subset \mathbb{R}
\]
が成立する。

\begin{remark}[自己共役性と実スペクトル]
自己共役作用素は量子力学における観測可能量に対応し、
そのスペクトルは測定可能な値の集合となる。
\end{remark}

%%%%%%%%%%%%%%%%%%%%%%%%%%%%%%%%%%%%%%%%%%%%%%%%%%%%%%%%%%%%%%
\subsection{スペクトル合併定理の詳細な証明}

有限群 $G$ に対して、正則表現は既約表現の直和に分解される。
すなわち、
\[
\lambda \cong \bigoplus_{\rho \in \widehat{G}} d_\rho \rho.
\]

したがって、
\[
\lambda(x) \cong \bigoplus_{\rho \in \widehat{G}} d_\rho \widehat{x}(\rho).
\]

作用素のスペクトルは直和に対して
\[
\sigma\left(\bigoplus_i A_i\right) = \bigcup_i \sigma(A_i)
\]
が成立するため、
\[
\sigma(x) = \bigcup_{\rho \in \widehat{G}} \sigma(\widehat{x}(\rho))
\]
が従う。

\begin{proof}[詳細な証明]
直和作用素 $A=\bigoplus_i A_i$ に対し、
\[
\lambda I-A = \bigoplus_i (\lambda I-A_i)
\]
である。可逆性は各ブロックの可逆性と同値であるため、
\[
\lambda I-A\ \text{不可逆} \iff \exists i,\ \lambda I-A_i\ \text{不可逆}.
\]
よってスペクトルは合併となる。
\end{proof}

%%%%%%%%%%%%%%%%%%%%%%%%%%%%%%%%%%%%%%%%%%%%%%%%%%%%%%%%%%%%%%
\subsection{スペクトル半径と群の幾何}

スペクトル半径は
\[
r(x) = \max_{\lambda \in \sigma(x)} |\lambda|
\]
で定義される。

また、Gelfand の公式により
\[
r(x) = \lim_{n \to \infty} \|x^n\|^{1/n}
\]
が成立する。

\begin{remark}[Gelfand 公式の意味]
この公式は、スペクトル半径が「長期的な成長率」を測る量であることを示す。
\end{remark}

この公式は、群の成長度やランダムウォークの性質と深く関係している。

例えば、$G$ が指数成長を持つ場合、ある種の元 $x$ に対して
\[
r(x) > \|x\|_2
\]
が成立することが知られている。

\begin{example}[ランダムウォークとスペクトル半径]
対称確率測度 $\mu$ に対し、
\[
M=\sum_g \mu(g)\lambda(g)
\]
のスペクトル半径はランダムウォークの遷移確率の減衰率を決定する。
\end{example}

%%%%%%%%%%%%%%%%%%%%%%%%%%%%%%%%%%%%%%%%%%%%%%%%%%%%%%%%%%%%%%
\section{von Neumann 環への拡張}

群環 $\mathbb{C}G$ をヒルベルト空間 $\ell^2(G)$ 上の作用素として実現したことで、
群環は単なる代数的対象ではなく、解析的対象としての側面を持つようになった。
ここでは、作用素ノルムによる完備化を通じて $C^*$-環を構成し、
さらに弱作用素位相による閉包を通じて von Neumann 環を構成する。

これらの構造は、群環の解析的性質を理解するために不可欠であり、
特にスペクトル半径、トレース、自己共役性、因子分解などの概念が自然に導入される。

%%%%%%%%%%%%%%%%%%%%%%%%%%%%%%%%%%%%%%%%%%%%%%%%%%%%%%%%%%%%%%
\subsection{作用素ノルムの完備化と reduced group algebra}

作用素ノルム
\[
\|x\| = \|\lambda(x)\|_{\mathrm{op}}
= \sup_{\|f\|=1} \|\lambda(x)f\|
\]
は、群環の元を作用素として扱う際の基本的なノルムである。

このノルムで群環を完備化すると、
\[
C^*_r(G) = \overline{\mathbb{C}G}^{\|\cdot\|}
\]
が得られる。これを **既約表現ノルム完備群環（reduced group $C^*$-algebra）** と呼ぶ。

この $C^*$-環は、群の構造を解析的に反映する極めて重要な対象であり、
特に以下の性質を持つ。

\paragraph{自己共役性}
\[
\|x^*x\| = \|x\|^2.
\]

\paragraph{スペクトル半径公式}
\[
r(x) = \lim_{n \to \infty} \|x^n\|^{1/n}.
\]

\paragraph{ユニタリ元の存在}
群の元 $g$ はユニタリ作用素 $\lambda(g)$ として実現されるため、
\[
\|\lambda(g)\| = 1.
\]

%%%%%%%%%%%%%%%%%%%%%%%%%%%%%%%%%%%%%%%%%%%%%%%%%%%%%%%%%%%%%%
\subsection{弱作用素位相による閉包と群 von Neumann 環}

$C^*$-環 $C^*_r(G)$ をさらに弱作用素位相で閉包すると、
\[
L(G) = \{\lambda(g) \mid g \in G\}''
\]
が得られる。これを **群 von Neumann 環** と呼ぶ。

ここで $S''$ は $S$ の両側可換子の可換子であり、
von Neumann の二重可換子定理により、弱作用素閉包と一致する。

\begin{definition}[群 von Neumann 環]
\[
L(G)=\lambda(G)''
\]
を群 von Neumann 環と呼ぶ。
\end{definition}

%%%%%%%%%%%%%%%%%%%%%%%%%%%%%%%%%%%%%%%%%%%%%%%%%%%%%%%%%%%%%%
\subsection{群 von Neumann 環のトレース構造}

群 von Neumann 環 $L(G)$ には自然なトレース
\[
\tau(x) = \langle \lambda(x)e_e, e_e \rangle
\]
が存在する。

$x = \sum_{g \in G} a_g g$ に対して、
\[
\tau(x) = a_e
\]
が成立する。

\begin{remark}[単位元係数の意味]
$a_e$ は「原点に戻る確率振幅」として解釈できる。
\end{remark}

%%%%%%%%%%%%%%%%%%%%%%%%%%%%%%%%%%%%%%%%%%%%%%%%%%%%%%%%%%%%%%
\subsection{トレースの性質}

トレース $\tau$ は以下の性質を満たす。

\paragraph{正値性}
\[
\tau(x^*x) \ge 0.
\]

\paragraph{正規性}
弱作用素連続である。

\paragraph{トレース性}
\[
\tau(xy) = \tau(yx).
\]

\paragraph{忠実性}
\[
\tau(x^*x) = 0 \Rightarrow x = 0.
\]

%%%%%%%%%%%%%%%%%%%%%%%%%%%%%%%%%%%%%%%%%%%%%%%%%%%%%%%%%%%%%%
\subsection{密度行列と量子状態の構成}

トレースを用いることで、群 von Neumann 環上の量子状態を構成できる。

正値作用素 $a \in L(G)$ に対して、
\[
\omega_a(x) = \tau(ax)
\]
は状態（positive linear functional）となる。

特に、$a$ がトレース 1 を満たす場合、
\[
\tau(a) = 1,
\]
$a$ は量子力学における密度行列と同じ役割を果たす。

%%%%%%%%%%%%%%%%%%%%%%%%%%%%%%%%%%%%%%%%%%%%%%%%%%%%%%%%%%%%%%
%%%%%%%%%%%%%%%%%%%%%%%%%%%%%%%%%%%%%%%%%%%%%%%%%%%%%%%%%%%%%%
\section{量子化と遷移作用素としての群環}

群環 $\mathbb{C}G$ の元を量子状態として解釈する視点は、
単なる比喩ではなく、数学的に厳密な構造を持つ。
特に、群環の積が「量子遷移作用素」として自然に理解できることは、
群環の解析的性質を量子力学的枠組みで説明する強力な方法となる。

量子化の観点から群環を見ると、群の離散的対称性が
量子状態の空間にそのまま反映される。
これは、量子ウォーク・量子情報・非可換確率論などの分野と
深く結びついている。

%%%%%%%%%%%%%%%%%%%%%%%%%%%%%%%%%%%%%%%%%%%%%%%%%%%%%%%%%%%%%%
\subsection{群環の元を量子状態として解釈する}

量子力学では、状態はヒルベルト空間のベクトルで表される。
群環の元
\[
x = \sum_{g \in G} a_g g
\]
は、係数 $a_g$ を「振幅」とみなすことで、
離散的な量子状態の重ね合わせとして解釈できる。

すなわち、$g$ は「状態」、$a_g$ はその状態の複素振幅である。

\begin{definition}[群環の量子状態]
$x=\sum_g a_g g$ が
\[
\sum_{g\in G}|a_g|^2=1
\]
を満たすとき、$x$ を量子状態とみなす。
\end{definition}

このとき、規格化条件
\[
\sum_{g \in G} |a_g|^2 = 1
\]
を課すことで、$x$ は量子状態として扱える。

\begin{remark}[離散量子系との同型性]
$G$ が有限群であれば、$\mathbb{C}G$ は有限次元量子系の Hilbert 空間と同型である。
無限群の場合は可分 Hilbert 空間となり、量子場の離散モデルに近い構造を持つ。
\end{remark}

%%%%%%%%%%%%%%%%%%%%%%%%%%%%%%%%%%%%%%%%%%%%%%%%%%%%%%%%%%%%%%
\subsection{群環の積と量子遷移振幅の合成}

群環の積
\[
xy = \sum_{k \in G} \left( \sum_{gh=k} a_g b_h \right) k
\]
は、量子状態の遷移振幅の合成と一致する。

すなわち、状態 $g$ から $h$ を経由して $k$ に到達する振幅は
\[
a_g b_h
\]
であり、すべての経路の振幅を足し合わせることで
\[
c_k = \sum_{gh=k} a_g b_h
\]
が得られる。

\begin{proposition}[群環積と量子遷移の同型性]
群環の積は、量子力学における遷移振幅の合成則と同型である。
\end{proposition}

\begin{proof}
量子遷移の合成則
\[
A(g\to k)=\sum_h A(g\to h)A(h\to k)
\]
と群環の積の係数
\[
c_k=\sum_{gh=k}a_g b_h
\]
が一致するため。
\end{proof}

これは、量子力学における「経路積分」の離散版と完全に同型である。

%%%%%%%%%%%%%%%%%%%%%%%%%%%%%%%%%%%%%%%%%%%%%%%%%%%%%%%%%%%%%%
\subsection{自己共役元と観測可能量}

$x = x^*$ が成立する場合、$\lambda(x)$ は自己共役作用素となり、
量子力学における「観測可能量」として解釈できる。

このとき、スペクトル $\sigma(x)$ は実数であり、
固有値は観測可能な値を表す。

特に、固有値分解
\[
x = \sum_{\lambda \in \sigma(x)} \lambda P_\lambda
\]
は、量子測定の結果を表す。

\begin{remark}[観測の確率構造]
射影 $P_\lambda$ は測定結果 $\lambda$ が得られる確率を決定する。
\end{remark}

%%%%%%%%%%%%%%%%%%%%%%%%%%%%%%%%%%%%%%%%%%%%%%%%%%%%%%%%%%%%%%
\section{非可換調和解析の深化}

群環の Fourier 解析は、単なる表現論の道具ではなく、
作用素環論・確率論・量子力学と深く結びついた強力な解析手法である。

%%%%%%%%%%%%%%%%%%%%%%%%%%%%%%%%%%%%%%%%%%%%%%%%%%%%%%%%%%%%%%
\subsection{Fourier 変換と作用素の分解}

群環の Fourier 変換
\[
\widehat{x}(\rho) = \sum_{g \in G} a_g \rho(g)
\]
は、作用素 $\lambda(x)$ の分解を与える。

実際、
\[
\lambda(x) \cong \bigoplus_{\rho \in \widehat{G}} d_\rho \widehat{x}(\rho)
\]
が成立する。

\begin{remark}[行列ブロック分解]
非可換 Fourier 変換は、作用素を既約表現ごとの行列に分解する操作である。
\end{remark}

%%%%%%%%%%%%%%%%%%%%%%%%%%%%%%%%%%%%%%%%%%%%%%%%%%%%%%%%%%%%%%
\subsection{Plancherel の等式とエネルギー分散}

群環の $\ell^2$ ノルム
\[
\|x\|_2^2 = \sum_{g \in G} |a_g|^2
\]
は、Fourier 変換を通じて
\[
\|x\|_2^2 = \frac{1}{|G|} \sum_{\rho \in \widehat{G}} d_\rho \|\widehat{x}(\rho)\|_{\mathrm{HS}}^2
\]
と表される。

これは、群環の元の「エネルギー」が既約表現の行列に分散していることを意味する。

\begin{remark}[エネルギーの分散]
可換群ではエネルギーは周波数成分に分散するが、
非可換群では行列ブロックに分散する。
\end{remark}

%%%%%%%%%%%%%%%%%%%%%%%%%%%%%%%%%%%%%%%%%%%%%%%%%%%%%%%%%%%%%%
\subsection{畳み込み作用と確率過程}

群環の畳み込み積は、確率論におけるランダムウォークと深く関係している。

確率測度 $\mu : G \to [0,1]$ に対して、
\[
\mu^{*n}(g) = \sum_{g_1 g_2 \cdots g_n = g} \mu(g_1)\mu(g_2)\cdots\mu(g_n)
\]
は $n$ ステップ後の分布を表す。

この構造は、群環の積
\[
x^n = x * x * \cdots * x
\]
と完全に一致する。

%%%%%%%%%%%%%%%%%%%%%%%%%%%%%%%%%%%%%%%%%%%%%%%%%%%%%%%%%%%%%%
\subsection{スペクトル半径とランダムウォーク}

ランダムウォークの遷移作用素 $T$ のスペクトル半径は、
群の幾何的性質と深く関係している。

特に、
\[
r(T) = \lim_{n \to \infty} \|T^n\|^{1/n}
\]
は、群の成長度やアーメナブル性と密接に結びついている。

例えば、$G$ が非アーメナブルである場合、
\[
r(T) < 1
\]
が成立することが知られている。

%%%%%%%%%%%%%%%%%%%%%%%%%%%%%%%%%%%%%%%%%%%%%%%%%%%%%%%%%%%%%%
\section{群 von Neumann 環の確率論的側面}

群 von Neumann 環 $L(G)$ は、非可換確率論の自然な舞台である。

%%%%%%%%%%%%%%%%%%%%%%%%%%%%%%%%%%%%%%%%%%%%%%%%%%%%%%%%%%%%%%
\subsection{トレースと期待値}

トレース $\tau$ は、確率論における期待値と同型の構造を持つ。

$x = \sum a_g g$ に対して、
\[
\tau(x) = a_e
\]
は「単位元に戻る確率振幅」を表す。

%%%%%%%%%%%%%%%%%%%%%%%%%%%%%%%%%%%%%%%%%%%%%%%%%%%%%%%%%%%%%%
\subsection{自由確率論との接続}

$G$ が自由群の場合、$L(G)$ は自由確率論の基本例となる。

特に、自由群の生成元 $u_1, \dots, u_n$ は自由独立性を満たし、
そのスペクトル分布は semicircular law に従う。

%%%%%%%%%%%%%%%%%%%%%%%%%%%%%%%%%%%%%%%%%%%%%%%%%%%%%%%%%%%%%%
\section{量子化のさらなる展開}

群環 $\mathbb{C}G$ を量子系として解釈する視点は、単なる比喩ではなく、
作用素環論と量子力学の構造が本質的に一致していることに基づく。

%%%%%%%%%%%%%%%%%%%%%%%%%%%%%%%%%%%%%%%%%%%%%%%%%%%%%%%%%%%%%%
\subsection{群環の元の「重ね合わせ」と量子状態の一致}

群環の元
\[
x = \sum_{g \in G} a_g g
\]
は、量子状態の重ね合わせと同型である。

規格化条件
\[
\|x\|_2^2 = \sum_{g \in G} |a_g|^2 = 1
\]
を課すことで、$x$ は量子状態として扱える。

%%%%%%%%%%%%%%%%%%%%%%%%%%%%%%%%%%%%%%%%%%%%%%%%%%%%%%%%%%%%%%
\subsection{群環の積と量子遷移作用素の一致}

群環の積
\[
xy = \sum_{k \in G} \left( \sum_{gh=k} a_g b_h \right) k
\]
は、量子遷移作用素の合成と同型である。

%%%%%%%%%%%%%%%%%%%%%%%%%%%%%%%%%%%%%%%%%%%%%%%%%%%%%%%%%%%%%%
\subsection{自己共役元と量子観測の一致}

$x = x^*$ のとき、$\lambda(x)$ は観測可能量となり、
スペクトル分解
\[
x = \sum_{\lambda \in \sigma(x)} \lambda P_\lambda
\]
は量子測定の構造を与える。

%%%%%%%%%%%%%%%%%%%%%%%%%%%%%%%%%%%%%%%%%%%%%%%%%%%%%%%%%%%%%%
\subsection{群 von Neumann 環と量子力学の一致}

群 von Neumann 環 $L(G)$ は量子力学の観測可能量の空間と同型であり、
トレース
\[
\tau(x)=\langle \lambda(x)e_e,e_e\rangle
\]
は期待値と一致する。

密度行列 $a$ に対し、
\[
\omega_a(x)=\tau(ax)
\]
は量子状態を表す。

%%%%%%%%%%%%%%%%%%%%%%%%%%%%%%%%%%%%%%%%%%%%%%%%%%%%%%%%%%%%%%
%%%%%%%%%%%%%%%%%%%%%%%%%%%%%%%%%%%%%%%%%%%%%%%%%%%%%%%%%%%%%%
\section{量子化のさらなる展開}

群環 $\mathbb{C}G$ を量子系として理解するためには、非可換代数としての構造と
ヒルベルト空間上の作用素としての構造がどのように整合するかを明確にする必要がある。
この節では、群環の元を量子状態、遷移作用素、観測可能量として扱うための
数学的基盤を、より厳密な観点から再構成する。

%%%%%%%%%%%%%%%%%%%%%%%%%%%%%%%%%%%%%%%%%%%%%%%%%%%%%%%%%%%%%%
\subsection{群環の元と量子状態の同型的構造}

群環 $\mathbb{C}G$ は有限支援関数空間と同一視でき、標準基底 $\{g\}_{g\in G}$ は
$\ell^2(G)$ の正規直交基底と一致する。したがって、群環の元
\[
x=\sum_{g\in G}a_g g
\]
は、係数列 $(a_g)$ が $\ell^2(G)$ に属する限り、量子状態と同型の構造を持つ。
規格化条件
\[
\sum_{g\in G}|a_g|^2=1
\]
が成立する場合、$x$ は離散量子系の純粋状態として扱うことができる。
このとき、内積
\[
\langle x,y\rangle=\sum_{g\in G}a_g\overline{b_g}
\]
は遷移振幅を与え、群の元が状態空間の正規直交基底として機能する。

%%%%%%%%%%%%%%%%%%%%%%%%%%%%%%%%%%%%%%%%%%%%%%%%%%%%%%%%%%%%%%
\subsection{群環の積と遷移作用素の構造}

群環の積は係数関数の畳み込みに対応し、
\[
xy=\sum_{k\in G}(a*b)(k)\,k,\qquad (a*b)(k)=\sum_{gh=k}a_g b_h
\]
で与えられる。この構造は、量子力学における遷移振幅の合成則と同型である。
$g$ から $h$ を経由して $k$ に至る遷移の振幅が $a_g b_h$ で与えられると解釈すると、
$k$ に至る全振幅はすべての経路の干渉和として表される。
非可換性は遷移の順序依存性として現れ、群の構造が遷移作用素の非可換性に直接反映される。

%%%%%%%%%%%%%%%%%%%%%%%%%%%%%%%%%%%%%%%%%%%%%%%%%%%%%%%%%%%%%%
\subsection{自己共役元と観測可能量}

群環の対合
\[
x^*=\sum_{g\in G}\overline{a_g}\,g^{-1}
\]
は、左正則表現 $\lambda$ を通じて自己共役作用素へ写る。
$x=x^*$ が成立する場合、$\lambda(x)$ は自己共役作用素となり、
そのスペクトルは実数に含まれる。
このとき、スペクトル分解
\[
\lambda(x)=\int_{\sigma(x)}\lambda\, dE(\lambda)
\]
が存在し、$E(\lambda)$ は観測値に対応する射影を与える。
量子測定の数学的枠組みはこのスペクトル分解により完全に記述されるため、
群環の自己共役元は観測可能量と同型の構造を持つ。

%%%%%%%%%%%%%%%%%%%%%%%%%%%%%%%%%%%%%%%%%%%%%%%%%%%%%%%%%%%%%%
\section{作用素環論との統合的視点}

群環を作用素として扱うことで、群の幾何的性質が作用素の解析的性質にどのように反映されるかを理解できる。
特に、スペクトル半径、アーメナブル性、自由確率論との関係は、
群の構造と作用素の構造が深く結びついていることを示す。

%%%%%%%%%%%%%%%%%%%%%%%%%%%%%%%%%%%%%%%%%%%%%%%%%%%%%%%%%%%%%%
\subsection{スペクトル半径と群の成長度}

スペクトル半径
\[
r(x)=\lim_{n\to\infty}\|x^n\|^{1/n}
\]
は、畳み込みの反復による広がりの速度を測る量である。
群が指数成長を持つ場合、畳み込みの反復は急速に広がるため、
適当な $x\in\mathbb{C}G$ に対して
\[
r(x)>\|x\|_2
\]
が成立することがある。
この現象は、群の成長度が作用素ノルムに直接影響することを示している。

%%%%%%%%%%%%%%%%%%%%%%%%%%%%%%%%%%%%%%%%%%%%%%%%%%%%%%%%%%%%%%
\subsection{アーメナブル性とスペクトル半径の等式}

有限生成群 $G$ の対称確率測度 $\mu$ に対し、遷移作用素
\[
T_\mu f=\mu*f
\]
のスペクトル半径は、アーメナブル性と密接に関係する。
Kesten の定理によれば、$G$ がアーメナブルであることと
\[
r(T_\mu)=1
\]
が同値であり、非アーメナブル群では $r(T_\mu)<1$ が成立する。
この等式は、群の幾何的性質が作用素のスペクトルに直接反映される典型例である。

%%%%%%%%%%%%%%%%%%%%%%%%%%%%%%%%%%%%%%%%%%%%%%%%%%%%%%%%%%%%%%
\subsection{自由群と自由確率論}

自由群 $F_n$ の群 von Neumann 環 $L(F_n)$ は自由確率論の基本例であり、
標準生成元の自己共役化が半円分布に従うことはよく知られている。
Voiculescu による自由確率論の枠組みでは、
自由独立性が群の還元語の構造と一致し、
$L(F_n)$ の解析的性質が自由確率論の中心的結果と自然に結びつく。

%%%%%%%%%%%%%%%%%%%%%%%%%%%%%%%%%%%%%%%%%%%%%%%%%%%%%%%%%%%%%%
\section{結論：群環の線形結合の数学的本質}

群環の線形結合
\[
x=\sum_g a_g g
\]
は、非可換代数としての構造、ヒルベルト空間上の作用素としての構造、
そして量子状態としての構造を同時に備えている。
これら三つの構造が矛盾なく統合されることにより、
群環は代数、解析、量子の三領域を結ぶ中心的対象として位置づけられる。

%%%%%%%%%%%%%%%%%%%%%%%%%%%%%%%%%%%%%%%%%%%%%%%%%%%%%%%%%%%%%%
\section{群環の解析的構造の総合}

群環の解析的構造を理解するためには、代数的側面と作用素論的側面を統合的に扱う必要がある。
群環 $\mathbb{C}G$ は有限支援関数空間としての線形構造を持つ一方で、
左正則表現を通じてヒルベルト空間 $\ell^2(G)$ 上の作用素として実現される。
この二重の構造が、スペクトル理論・作用素ノルム・トレースといった解析的概念を導入する基盤となる。

%%%%%%%%%%%%%%%%%%%%%%%%%%%%%%%%%%%%%%%%%%%%%%%%%%%%%%%%%%%%%%
\subsection{スペクトル理論の統合的役割}

群環の元 $x=\sum_{g\in G}a_g g$ に対し、左正則表現 $\lambda$ を通じて得られる作用素
$\lambda(x)$ のスペクトル
\[
\sigma(x)=\{\lambda\in\mathbb{C}\mid \lambda I-\lambda(x)\ \text{が可逆でない}\}
\]
は、$x$ の解析的性質を決定する中心的対象である。
スペクトルは常に非空であり、コンパクトであり、作用素ノルムによって有界である。
これらの性質は一般の有界作用素に対する標準的結果であるが、
群環の場合には表現論的分解と結びつくことで、より具体的な構造が得られる。

有限群の場合、正則表現は既約表現の直和に分解されるため、
\[
\lambda(x)\cong\bigoplus_{\rho\in\widehat{G}} d_\rho\,\widehat{x}(\rho)
\]
が成立する。
このとき、作用素のスペクトルが直和に対して合併として振る舞うことから、
\[
\sigma(x)=\bigcup_{\rho\in\widehat{G}}\sigma(\widehat{x}(\rho))
\]
が従う。
この等式は、群環の解析が既約表現の行列解析へと完全に還元されることを意味し、
群環の解析的構造が表現論的構造と密接に結びついていることを示す。

%%%%%%%%%%%%%%%%%%%%%%%%%%%%%%%%%%%%%%%%%%%%%%%%%%%%%%%%%%%%%%
\subsection{作用素ノルムとスペクトル半径の関係}

作用素ノルム
\[
\|x\|=\sup_{\|f\|=1}\|\lambda(x)f\|
\]
は、群環の元を作用素として扱う際の基本的量である。
このノルムは群の幾何的性質を反映し、特に畳み込みの反復に対する成長率と密接に関係する。

Gelfand の公式
\[
r(x)=\lim_{n\to\infty}\|x^n\|^{1/n}
\]
は、スペクトル半径が作用素ノルムの漸近的挙動によって決定されることを示す。
この公式は、群の成長度やランダムウォークの漸近挙動と深く結びついており、
群の幾何的性質が作用素の解析的性質に直接反映される典型例である。
特に、指数成長を持つ群では、畳み込みの反復が急速に広がるため、
スペクトル半径が $\ell^2$ ノルムを上回る現象が生じる。

%%%%%%%%%%%%%%%%%%%%%%%%%%%%%%%%%%%%%%%%%%%%%%%%%%%%%%%%%%%%%%
\subsection{群 von Neumann 環のトレース構造}

群 von Neumann 環 $L(G)$ は、左正則表現の像の弱作用素閉包として定義される。
この環には自然なトレース
\[
\tau(x)=\langle \lambda(x)e_e,e_e\rangle=a_e
\]
が存在し、これは群環の元の単位元成分を抽出する線形汎関数である。
このトレースは正値性・正規性・トレース性・忠実性を満たし、
量子力学における期待値と同型の構造を持つ。

正値作用素 $a\in L(G)$ に対し、
\[
\omega_a(x)=\tau(ax)
\]
と定めると、これは $L(G)$ 上の状態を与える。
この構造は、群 von Neumann 環が量子力学の観測可能量の空間と同型の構造を持つことを示している。

%%%%%%%%%%%%%%%%%%%%%%%%%%%%%%%%%%%%%%%%%%%%%%%%%%%%%%%%%%%%%%
\section{量子化の総合的理解}

群環の元
\[
x=\sum_{g\in G}a_g g
\]
は、係数列 $(a_g)$ が $\ell^2(G)$ に属する限り、量子状態の重ね合わせと同型の構造を持つ。
規格化条件
\[
\sum_{g\in G}|a_g|^2=1
\]
を課すことで、$x$ は離散量子系の純粋状態として扱うことができる。
このとき、内積
\[
\langle x,y\rangle=\sum_{g\in G}a_g\overline{b_g}
\]
は遷移振幅を与え、群の元が状態空間の正規直交基底として機能する。

群環の積
\[
xy=\sum_{k\in G}\left(\sum_{gh=k}a_g b_h\right)k
\]
は、量子遷移振幅の合成則と一致する。
$g$ から $h$ を経由して $k$ に至る振幅が $a_g b_h$ で与えられると解釈すると、
$k$ に至る全振幅はすべての経路の干渉和として表される。
この構造は、群環が量子力学の形式を自然に包含していることを示している。

%%%%%%%%%%%%%%%%%%%%%%%%%%%%%%%%%%%%%%%%%%%%%%%%%%%%%%%%%%%%%%
\section{今後の展望}

群環の解析的研究は、有限群にとどまらず無限群へと拡張される。
無限群においてはアーメナブル性が中心的役割を果たし、
特に Kesten の定理
\[
G\ \text{がアーメナブル} \iff r(T)=1
\]
は、群の幾何と作用素のスペクトル構造が密接に結びついていることを示す。
アーメナブル性の有無は、ランダムウォークの漸近挙動や作用素ノルムの成長に直接影響する。

群環に零因子が存在するかどうかという問題は、Kaplansky の零因子予想として知られ、
現代数学における重要な未解決問題である。
特に自由群に対して零因子が存在しないかどうかは深い問題であり、
群環の代数的構造と解析的構造の双方を巻き込む。

自由群の群 von Neumann 環 $L(F_n)$ は自由確率論の基本例であり、
自由独立性や半円分布といった概念が群環の解析と自然に結びつく。
この接続は、非可換確率論と群環の解析を統合する新たな方向性を提供する。

さらに、群環の構造は量子群へと自然に一般化される。
量子群の作用素環は非可換幾何学の中心的対象であり、
群環の解析的手法はそのまま量子群の研究へと拡張される。
この一般化は、非可換空間の幾何学的理解を深める上で重要である。

%%%%%%%%%%%%%%%%%%%%%%%%%%%%%%%%%%%%%%%%%%%%%%%%%%%%%%%%%%%%%%
\section{総括}

群環の線形結合
\[
x=\sum_{g\in G}a_g g
\]
は、代数的構造、解析的構造、量子的構造を同時に備えた極めて豊かな数学的対象である。
本稿では、群環を作用素として扱うことで、スペクトル理論、作用素ノルム、トレース、
量子状態といった概念が自然に導入されることを示し、
群環の解析的理解がの有用性を示した。
このように、群環の線形結合は単なる形式和ではなく、代数的・非代数的解析対象であることがわかる。

%%%%%%%%%%%%%%%%%%%%%%%%%%%%%%%%%%%%%%%%%%%%%%%%%%%%%%%%%%%%%%
\pagebreak
%%%%%%%%%%%%%%%%%%%%%%%%%%%%%%%%%%%%%%%%%%%%%%%%
\begin{thebibliography}{99}

\bibitem{dixon}
J. D. Dixon and B. Mortimer, \textit{Permutation Groups}, Graduate Texts in Mathematics, vol. 163, Springer-Verlag, New York, 1996.

\bibitem{passman}
D. S. Passman, \textit{The Algebraic Structure of Group Rings}, Wiley-Interscience, New York, 1977.

\bibitem{takesaki1}
M. Takesaki, \textit{Theory of Operator Algebras I}, Encyclopaedia of Mathematical Sciences, vol. 124, Springer-Verlag, Berlin, 2001.

\bibitem{takesaki2}
M. Takesaki, \textit{Theory of Operator Algebras II}, Encyclopaedia of Mathematical Sciences, vol. 125, Springer-Verlag, Berlin, 2003.

\bibitem{takesaki3}
M. Takesaki, \textit{Theory of Operator Algebras III}, Encyclopaedia of Mathematical Sciences, vol. 127, Springer-Verlag, Berlin, 2003.

\bibitem{connes}
A. Connes, \textit{Noncommutative Geometry}, Academic Press, San Diego, 1994.

\bibitem{blackadar}
B. Blackadar, \textit{Operator Algebras: Theory of C*-Algebras and von Neumann Algebras}, Encyclopaedia of Mathematical Sciences, vol. 122, Springer-Verlag, Berlin, 2006.

\bibitem{kesten1}
H. Kesten, Full Banach mean values on countable groups, \textit{Math. Scand.}, \textbf{7} (1959), 146--156.

\bibitem{kesten2}
H. Kesten, Symmetric random walks on groups, \textit{Trans. Amer. Math. Soc.}, \textbf{92} (1959), 336--354.

\bibitem{woess}
W. Woess, \textit{Random Walks on Infinite Graphs and Groups}, Cambridge Tracts in Mathematics, vol. 138, Cambridge University Press, 2000.

\bibitem{dirac}
P. A. M. Dirac, \textit{The Principles of Quantum Mechanics}, 4th Edition, Oxford University Press, 1958.

\bibitem{feynman}
R. P. Feynman and A. R. Hibbs, \textit{Quantum Mechanics and Path Integrals}, McGraw-Hill, New York, 1965.

\bibitem{kadison1}
R. V. Kadison and J. R. Ringrose, \textit{Fundamentals of the Theory of Operator Algebras. Vol. I: Elementary Theory}, Academic Press, 1983.

\bibitem{kadison2}
R. V. Kadison and J. R. Ringrose, \textit{Fundamentals of the Theory of Operator Algebras. Vol. II: Advanced Theory}, Academic Press, 1986.

\bibitem{serre}
J.-P. Serre, \textit{Linear Representations of Finite Groups}, Graduate Texts in Mathematics, vol. 42, Springer-Verlag, 1977.

\bibitem{pedersen}
G. K. Pedersen, \textit{C*-Algebras and their Automorphism Groups}, Academic Press, London, 1979.

\bibitem{murphy}
G. J. Murphy, \textit{C*-Algebras and Operator Theory}, Academic Press, 1990.

\bibitem{davidson}
K. R. Davidson, \textit{C*-Algebras by Example}, Fields Institute Monographs, American Mathematical Society, 1996.

\bibitem{weyl}
H. Weyl, \textit{The Theory of Groups and Quantum Mechanics}, Dover Publications, 1950.

\bibitem{neumann}
J. von Neumann, \textit{Mathematical Foundations of Quantum Mechanics}, Princeton University Press, 1955.

\bibitem{arveson}
W. Arveson, \textit{An Invitation to C*-Algebras}, Graduate Texts in Mathematics, vol. 39, Springer-Verlag, 1976.

\bibitem{brown}
N. P. Brown and N. Ozawa, \textit{C*-Algebras and Finite-Dimensional Approximations}, Graduate Studies in Mathematics, vol. 88, American Mathematical Society, 2008.

\bibitem{sakai}
S. Sakai, \textit{C*-Algebras and W*-Algebras}, Classics in Mathematics, Springer-Verlag, 1971.

\bibitem{bourbaki}
N. Bourbaki, \textit{Spectral Theory}, Elements of Mathematics, Springer-Verlag, 2013.

\bibitem{halmos}
P. R. Halmos, \textit{A Hilbert Space Problem Book}, 2nd Edition, Graduate Texts in Mathematics, vol. 19, Springer-Verlag, 1982.

\end{thebibliography}

\end{document}